\section{Power-map-free subsets of symmetric groups: \texorpdfstring{\seq{A007234}}{A007234}}
\label{sec:a007234}

Let \(F_n\) be the maximum cardinality of a subset \(X\subseteq S_n\) such
that
\begin{equation}
 \sigma\in X\quad\Longrightarrow\quad \sigma^2\notin X.
 \label{eq:a007234-condition}
\end{equation}
Equivalently, \(F_n\) is the independence number of the functional graph of
the squaring map, with a looped vertex forbidden.  Bouchard and Yeh studied
this general \(f\)-free-set problem and reported the first values
\citep{BouchardYeh1992}.  The conjugacy-type recurrence below corrects the last
two of those values and extends the sequence to every \(n\).

\subsection{Cycle types and square-root multiplicities}

Write a cycle type as \(\lambda=1^{m_1}2^{m_2}\cdots n^{m_n}\) and put
\begin{equation}
 z_\lambda=\prod_{j=1}^{n}j^{m_j}m_j!.
 \label{eq:a007234-z}
\end{equation}
The conjugacy class has size \(n!/z_\lambda\).  Let \(q(\mu)\) be the cycle
type obtained by squaring a permutation of type \(\mu\): an odd \(j\)-cycle
remains a \(j\)-cycle, while an even \(2j\)-cycle becomes two \(j\)-cycles.

\begin{lemma}[Uniform root fibers]
\label{lem:a007234-roots}
If \(q(\mu)=\lambda\), then every fixed permutation of type \(\lambda\) has
exactly
\begin{equation}
 r(\lambda,\mu)=\frac{z_\lambda}{z_\mu}
 \label{eq:a007234-roots}
\end{equation}
square roots of type \(\mu\).
\end{lemma}

\begin{proof}
The squaring map sends the conjugacy class of \(\mu\) equivariantly into the
class of \(\lambda\).  Its image is nonempty and invariant under conjugation;
hence it is the whole target class.  Equivariance also makes every fiber have
the same size.  Dividing the source-class size \(n!/z_\mu\) by the target-class
size \(n!/z_\lambda\) gives \cref{eq:a007234-roots}.
\end{proof}

\subsection{Reverse trees}

For a type \(\lambda\) containing an even part, let \(E_\lambda\) and
\(I_\lambda\) be the maximum sizes of an independent set in the reverse tree
below one permutation of type \(\lambda\), conditional on excluding or
including its root.  These values depend only on the type: conjugation carrying
one root permutation to another is an isomorphism of their rooted reverse
functional graphs.  Distinct square roots have disjoint reverse subtrees,
because every permutation has a unique square.  Therefore
\begin{align}
 E_\lambda&=\sum_{q(\mu)=\lambda}
 r(\lambda,\mu)\max\{E_\mu,I_\mu\},
 \label{eq:a007234-E}\\
 I_\lambda&=1+\sum_{q(\mu)=\lambda}
 r(\lambda,\mu)E_\mu.
 \label{eq:a007234-I}
\end{align}
These equations are acyclic.  Let
\[
 v(\lambda)=\max\{v_2(j):m_j(\lambda)>0\}.
\]
If \(\lambda\) has an even part and \(q(\mu)=\lambda\), then
\(v(\mu)=v(\lambda)+1\): a target cycle having maximal 2-adic valuation can
only arise from a cycle twice as long, and no source cycle can have larger
valuation without producing a larger target valuation.  Processing types in
descending \(v\) therefore determines all nonperiodic states.  The reverse
trees are finite because cycle lengths cannot be doubled indefinitely inside
\(S_n\).

If all parts of \(\lambda\) are odd, the same formulae determine the tree
attached to a periodic cycle vertex after omitting the unique same-type root.
Indeed, \(q(\lambda)=\lambda\) and
\(r(\lambda,\lambda)=1\); that root is the predecessor on the directed cycle,
not a child in the attached tree.

\subsection{Periodic conjugacy classes}

For an all-odd type \(\lambda\), put
\begin{equation}
 o_\lambda=\lcm\{j:m_j(\lambda)>0\},
 \qquad
 L_\lambda=\ord_{o_\lambda}(2),
 \label{eq:a007234-period}
\end{equation}
with \(L_\lambda=1\) when \(o_\lambda=1\).  Every permutation in the class has
order \(o_\lambda\), so its squaring orbit has length \(L_\lambda\).  The
class splits into
\begin{equation}
 N_\lambda=\frac{n!/z_\lambda}{L_\lambda}
 \label{eq:a007234-numbercycles}
\end{equation}
directed cycles.

Let \(g_\lambda=\max\{0,I_\lambda-E_\lambda\}\).  Excluding every cycle root
contributes \(L_\lambda E_\lambda\).  Including a nonadjacent set of cycle
roots gains \(g_\lambda\) per selected root, so a cycle of length greater than
one contributes
\begin{equation}
 L_\lambda E_\lambda+
 \left\lfloor\frac{L_\lambda}{2}\right\rfloor g_\lambda.
 \label{eq:a007234-cycleweight}
\end{equation}
When \(L_\lambda=1\), the root has a loop and is forbidden, so the contribution
is \(E_\lambda\).

\begin{theorem}[Conjugacy-type recurrence]
\label{thm:a007234}
Equations
\cref{eq:a007234-roots,eq:a007234-E,eq:a007234-I,eq:a007234-period,eq:a007234-numbercycles,eq:a007234-cycleweight}
form an exact finite recurrence for \(F_n\).  Explicitly,
\begin{equation}
 F_n=\sum_{\lambda\text{ all odd}}N_\lambda
 \begin{cases}
 E_\lambda,&L_\lambda=1,\\[2mm]
 L_\lambda E_\lambda+
 \lfloor L_\lambda/2\rfloor g_\lambda,&L_\lambda>1.
 \end{cases}
 \label{eq:a007234-main}
\end{equation}
\end{theorem}

\begin{proof}
Every finite functional graph is a disjoint union of directed cycles with
rooted in-trees attached.  The reverse-tree states
\cref{eq:a007234-E,eq:a007234-I} are the usual include/exclude dynamic
program, with root multiplicities supplied by
\cref{lem:a007234-roots}.  The valuation ordering proves that these states are
computed without circular dependence away from all-odd types.  On an all-odd
class the squaring map is a permutation with cycle length
\(L_\lambda\), and the same-type predecessor has been removed from the
attached tree.  The maximum-weight independent set on each directed cycle is
\cref{eq:a007234-cycleweight}, with the loop exception at length one.
Summing all components yields \cref{eq:a007234-main}.
\end{proof}

The recurrence begins, with offset one,
\[
\begin{aligned}
0,&\ 1,\ 4,\ 16,\ 72,\ 522,\ 3390,\ 29409,\ 267561,\\
&2820600,\ 30658050,\ 377859960,\ 4866471720,\\
&70224099120,\ 1052687890800,\ 17121988170000,\ldots.
\end{aligned}
\]
Thus the values formerly reported as \(3642\) and \(30753\) at \(n=7,8\) are
corrected to
\begin{equation}
 \boxed{F_7=3390,\qquad F_8=29409.}
 \label{eq:a007234-corrections}
\end{equation}

The recurrence agrees through \(n=50\) with its stored exact table.  Complete
functional graphs through \(S_{11}\), two independent Python component
algorithms, and an exhaustive audit of every square-root fiber through
\(S_7\) give structurally independent checks.  Explicit maximum-set rank lists
for \(n=7,8,9\) certify the lower-bound side directly.

\subsection{Every fixed power map}

The argument extends without a change of principle to
\(P_d:S_n\to S_n\), \(P_d(\sigma)=\sigma^d\), for any fixed \(d\ge2\).  A
cycle of length \(j\) becomes \(\gcd(j,d)\) cycles of length
\(j/\gcd(j,d)\).  Let \(q_d(\mu)\) be the induced type map.  Equivariance gives
\begin{equation}
 r_d(\lambda,\mu)=\frac{z_\lambda}{z_\mu}
 \qquad(q_d(\mu)=\lambda).
 \label{eq:a007234-general-roots}
\end{equation}

Let \(\mathcal P\) be the set of prime divisors of \(d\), and define
\begin{equation}
 h_d(\lambda)=\sum_{p\in\mathcal P}
   \max\{v_p(j):m_j(\lambda)>0\}.
 \label{eq:a007234-height}
\end{equation}
For a source cycle of length \(\ell\) and its target length
\(j=\ell/\gcd(\ell,d)\),
\begin{equation}
 v_p(j)=\max\{v_p(\ell)-v_p(d),0\}.
 \label{eq:a007234-valuation}
\end{equation}
If a target type has a part not coprime to \(d\), then every reverse type has
strictly larger \(h_d\): for at least one prime \(p\mid d\), a positive target
valuation forces the corresponding source valuation to increase by
\(v_p(d)\), while the maxima for other primes cannot decrease.  The reverse
type graph is therefore acyclic away from types whose parts are all coprime to
\(d\).

For such a periodic type, put
\begin{equation}
 o_\lambda=\lcm\{j:m_j(\lambda)>0\},
 \qquad
 L_\lambda=\ord_{o_\lambda}(d).
 \label{eq:a007234-general-period}
\end{equation}
The class splits into power-map cycles of length \(L_\lambda\); omitting the
unique same-type predecessor and applying the same weighted-cycle formula
computes the largest \(X\subseteq S_n\) satisfying
\[
 \sigma\in X\Longrightarrow \sigma^d\notin X.
\]
This gives an exact all-\(n\) recurrence for every fixed \(d\).  The generic
implementation specializes to \cref{thm:a007234} for \(d=2\) and agrees with
direct functional graphs for \(d=2,3,4,6\) through \(S_7\).
